\section{The Role of the Cabinet}

\subsection{Composition of the Cabinet}

The KGB represents the Soviet and allied intelligence, political, and military figures with an interest in the development of Latin America during the Cold War. The cabinet is centred on the Soviet Union, particularly its political leadership, foreign intelligence apparatus, and military establishment, while also incorporating key Cuban intelligence and political figures with direct knowledge of revolutionary movements in the region.

Its members therefore do not represent a single institution or possess identical priorities. Soviet political leaders are concerned with the broader international position of the USSR, while KGB officials focus more directly on intelligence, foreign networks, and the activities of political movements abroad. Military figures may approach the situation from a strategic and security perspective, while Cuban representatives bring their own regional experience and interests.

The KGB should be understood as a collection of actors operating within the broader Soviet sphere of influence, rather than as a single unified intelligence structure. Their interests may overlap, but their access to information, responsibilities, and perceptions of the situation can differ considerably.

\subsection{Objectives}

The KGB’s primary objective is to protect and advance Soviet interests within the wider Cold War. Latin America represents an opportunity to challenge American influence and strengthen the international position of the Soviet Union, but developments in the region must also be assessed according to their potential consequences for Soviet foreign policy and relations with Washington.

A central concern is the collection, verification, and interpretation of information regarding governments, revolutionary organizations, intelligence services, and foreign involvement in Latin America. The cabinet must determine what information is reliable, what remains uncertain, and which developments may have broader strategic significance.

The Soviet Union and its partners have an interest in the development of left-wing and revolutionary movements throughout Latin America. However, these organizations differ considerably in ideology, objectives, capabilities, and relationships with Moscow. The cabinet must therefore assess where political, intelligence, or material support could advance Soviet interests and where involvement could create unnecessary risks.

\subsection{Relationship to other Cabinets}

\subsection{Cabinet A - Regional Council}

The Regional Council constitutes the principal political authorities with which Soviet and Cuban actors must contend in the region. Their different political systems and relationships with the United States make them neither a completely uniform adversary nor an entirely predictable one. Understanding their internal divisions, political vulnerabilities, and relationships with one another may be as important as understanding their opposition to revolutionary movements.

\subsection{Cabinet B - CIA}

The United States represents the principal strategic competitor of the Soviet Union in Latin America. The CIA's intelligence activities and relationships with regional governments may provide important information about American intentions and capabilities. At the same time, American involvement creates opportunities for the Soviet Union to identify weaknesses in U.S. policy and potentially expand its own influence.

\subsection{Cabinet C - Revolutionaries}

Revolutionary organizations are potential partners as well as independent political actors. Their ideological diversity and varying degrees of commitment to Soviet interests mean that their objectives cannot necessarily be assumed to coincide with those of Moscow. Relations with these movements may therefore involve political support, intelligence cooperation, competition for influence, or strategic caution.
